\section*{Adaptaciones a CRISP-DM}
\addcontentsline{toc}{section}{Adaptaciones a CRISP-DM}

En esta fase hicimos los siguientes ajustes a CRISP-DM:

\begin{itemize}
  \item Juntamos los cuatro entregables de la fase 2 (reporte de recolección inicial, de descripción, de exploración y de calidad de los datos) en este documento, en secciones separadas.
  \item Ordenamos el reporte alrededor del BCS, que es el único objetivo del proyecto.
  \item Los datos de ordeño se quedan como contexto, porque solo 12 de las 69 vacas con BCS tienen datos de ordeño. Se resumen dentro de este mismo documento.
  \item Analizamos los datos de BCS y las fotografías en el notebook de Python \texttt{Reporte de Datos BCS fotos.ipynb}, que está en la ruta \texttt{src/notebooks/} del repositorio.
  \item Los riesgos que encontramos en los datos se evalúan con probabilidad, impacto y magnitud, como en el plan de riesgos de la fase 1. La diferencia es que aquí la probabilidad usa una escala de 1 a 5.
\end{itemize}
